% Propietario conservado: /Users/ruben/Documents/ChatGPT/jueces y controles/REVISION_ARGUMENTAL_APERTURAS_CIERRES_20260915/delivery/ARTICLE_V_ES_ARGUMENTAL_20260915/manuscrito/sections/40_exclusion_pauli.tex
% Recuperación documental para el artículo V; RESULTADO_RECUPERADO.
% Propietario: /Users/ruben/Documents/New project/output/REVISION_EDITORIAL_HMT_MD_20260905/01_FUENTES/INTEGRAL/tree/New project/output/ENSAMBLADO_CAUSAL_RESTAURADO_HMT_MD_20260905/fuente/manuscrito/integracion_83/parte_iv/body/B015_ch56_pauli_completacion_fermionica_7549b83967f5_04d_ocupacion_estadistica.tex
% Líneas de origen: 1--48.
% REV02: aplicación estadística; prueba común de idempotencia reunida en 20.
% Recupera además los recuentos y ocupaciones antes anticipados en 20.
% Correspondencia editorial: gestion/CONSOLIDACION_PAULI_FOCK_REV02.json.
\section{Aplicación estadística de los espectros de ocupación}
\label{v:v:rec:chap:ocupacion-estadistica}
\index{Pauli, principio de exclusión de}\index{CAR}
\index{Fermi--Dirac}\index{Bose--Einstein}

Las realizaciones nula y material proporcionan el soporte geométrico de la
ocupación. La completación graduada de ese soporte determina sus espacios
multiparticulares y sus operadores; el
teorema~\ref{v:v:rec:fock_gibbs:statement:01} establece su transporte por
refinamiento. Este tramo utiliza esos resultados para caracterizar las
configuraciones admisibles, contar estados y preparar la evaluación térmica.
La construcción algebraica, el recuento finito y el estado de equilibrio
conservan así sus funciones y premisas respectivas.

\subsection{Restricción espectral y configuraciones admisibles}
\label{v:v:rec:pauli:section:02}

Para un modo \(i\), escribimos \(n_i=N_i=a_i^\dagger a_i\). El
teorema~\ref{v:v:rec:thm:principio-exclusion-pauli}, con su prueba completa en
\eqref{v:v:pauli:eq:prueba-idempotencia}, establece
\(n_i^2=n_i\) y \(\Spec(n_i)\subseteq\{0,1\}\). El cálculo utiliza
únicamente CAR y la relación de adjunción, por lo que se aplica asimismo
a cualquier representación de esas relaciones. La nilpotencia
\((a_i^\dagger)^2=0\) impide la ocupación fermiónica doble del mismo modo.
Por consiguiente, una configuración fermiónica de \(g\) modos es una
familia \((n_1,\ldots,n_g)\in\{0,1\}^g\); en el sector bosónico, las
ocupaciones pertenecen a \(\mathbb N_0^g\).

La representación fermiónica utiliza el álgebra CAR construida desde la
completación exterior. La identificación relativista de
espín--estadística conserva, además, la compatibilidad de la representación
de rotaciones con la paridad y las condiciones de localidad expuestas en
las secciones~\ref{v:v:pauli:car-construccion} y
\ref{v:v:rec:completaciones:section:05}. La periodicidad \(4\pi\) y la
exclusión quedan coordinadas por el mismo carácter de hoja, con esas
condiciones de representación explícitas.

\subsection{Dimensiones de los sectores de ocupación fija}
\label{v:v:rec:completaciones:section:07}

Si hay \(g\) modos y \(N\) partículas, las dimensiones de los sectores son
\begin{equation}
 \Omega_F(g,N)=\binom gN,
 \qquad
 \Omega_B(g,N)=\binom{g+N-1}{N}.
 \label{v:v:ocupacion:eq:dimensiones-sectores}
\end{equation}
La primera fórmula cuenta subconjuntos de modos, con dimensión nula si
\(N>g\); la segunda, composiciones débiles de \(N\) en \(g\) partes.
Estos conteos son consecuencias exactas de las dos completaciones y se
obtienen antes de introducir el estado térmico. La
proposición~\ref{v:v:rec:prop:cardinalidades-microcanonicas} aplica los dos
recuentos a la pantalla espectral finita, conservando sus niveles,
multiplicidades y restricciones de energía y población.

\subsection{Lectura térmica de las completaciones graduadas}
\label{v:v:ocupacion:lectura-termica}

Con un Hamiltoniano diagonal \(H\), energía modal \(\epsilon\), potencial
químico \(\mu\) y parámetro térmico \(\beta>0\), las ocupaciones medias
de los dos sectores son
\begin{equation}
 \bar n_{\rm par}(\epsilon)
 =\frac{1}{e^{\beta(\epsilon-\mu)}-1},
 \qquad
 \bar n_{\rm impar}(\epsilon)
 =\frac{1}{e^{\beta(\epsilon-\mu)}+1}.
 \label{v:v:ocupacion:eq:leyes-por-modo}
\end{equation}
En el sector bosónico se requiere \(\mu<\inf\operatorname{Spec}H\).
Estas expresiones son la especialización por modo del
teorema~\ref{v:v:rec:thm:factorizacion-gran-canonica-rev3}, cuya prueba
construye la traza gran canónica a partir de las ocupaciones admisibles.
Las fórmulas de Bose--Einstein y Fermi--Dirac son, por tanto, las
evaluaciones termodinámicas de las completaciones simétrica y exterior
seleccionadas por el transporte. La misma sección obtiene sus expresiones
mediante la variación de la entropía bajo las restricciones de energía y
número; el teorema~\ref{v:v:rec:thm:limite-clasico-fd-be-rev2} demuestra el
límite clásico común con una cota explícita del error.
